\section{Introduction}

\subsection{Another incarnation of the Satake category}

Let $G$ be a connected reductive algebraic group over an algebraically closed field $\jepPubBbb{F}$ of positive characteristic, and let $\Bbbk$ be either a finite field of characteristic $\ell \neq \mathrm{char}(\jepPubBbb{F})$, or a finite extension of $\jepPubBbb{Q}_\ell$, or the ring of integers of such an extension. If $\jepPubScript{K}:=\jepPubBbb{F}( \hspace{-0.5pt} (z) \hspace{-0.5pt} )$ and $\jepPubScript{O}:= \jepPubBbb{F}[ \hspace{-0.5pt} [z] \hspace{-0.5pt} ]$, the \emph{Satake category} is the category
\[
 \mathrm{Perv}_{G_{\jepPubScript{O}}}(\mathrm{Gr},\Bbbk)
\]
of $G_{\jepPubScript{O}}$-equivariant (\'etale) $\Bbbk$-perverse sheaves on the affine Grassmannian
\[
 \mathrm{Gr}:=G_{\jepPubScript{K}}/G_{\jepPubScript{O}}
\]
of $G$. This category is a fundamental object in Geometric Representation Theory through its appearance in the \emph{geometric Satake equivalence}, which claims that this category admits a natural convolution product $(-) \star^{G_{\jepPubScript{O}}} (-)$, which endows it with a monoidal structure, and that there exists an equivalence of monoidal categories
\begin{equation}
\label{jep104local006}
 \jepPubEuler{S} : (\mathrm{Perv}_{G_{\jepPubScript{O}}}(\mathrm{Gr},\Bbbk), \star^{G_{\jepPubScript{O}}}) \xrightarrow{\sim} (\mathrm{Rep}(G^\vee_\Bbbk), \otimes).
\end{equation}
Here the right-hand side is the category of algebraic representations of the split reductive $\Bbbk$-group scheme which is Langlands dual to $G$ on finitely generated $\Bbbk$-modules; see~\cite{jep104ref36} for the original proof of this equivalence in full generality, and~\cite{jep104ref11} for a more detailed exposition. (In these references, what is explicitly treated is the analogous equivalence for a complex group $G$, in which case $\Bbbk$ can be any Noetherian commutative ring of finite global dimension. The \'etale setting is similar; see~\cite[\S 14]{jep104ref36} and~\cite[\S 1.1.4]{jep104ref11} for a few comments.)

This category already has another incarnation since (as proved by Mirkovi\'c--Vilonen) the forgetful functor
\[
 \mathrm{Perv}_{G_{\jepPubScript{O}}}(\mathrm{Gr},\Bbbk) \to \mathrm{Perv}_{(G_{\jepPubScript{O}})}(\mathrm{Gr},\Bbbk)
\]
from the Satake category to the category of perverse sheaves on $\mathrm{Gr}$ which are constructible with respect to the stratification by $G_{\jepPubScript{O}}$-orbits is an equivalence of categories. 

The first main result of the present paper provides a third incarnation of this category, as a category $\mathrm{Perv}_{\jepPubEuler{IW}}(\mathrm{Gr},\Bbbk)$ of Iwahori--Whittaker\footnote{This terminology is taken from~\cite{jep104ref10}. In~\cite{jep104ref09}, the term ``baby Whittaker'' is used for the same construction.} perverse sheaves on $\mathrm{Gr}$. More precisely we prove that a natural functor
\begin{equation}
\label{jep104local009}
 \mathrm{Perv}_{G_{\jepPubScript{O}}}(\mathrm{Gr},\Bbbk) \to \mathrm{Perv}_{\jepPubEuler{IW}}(\mathrm{Gr},\Bbbk)
\end{equation}
is an equivalence of categories, see Theorem~\ref{jep104local038}. This result is useful because computations in $\mathrm{Perv}_{\jepPubEuler{IW}}(\mathrm{Gr},\Bbbk)$ are much easier than in the categories $\mathrm{Perv}_{G_{\jepPubScript{O}}}(\mathrm{Gr},\Bbbk)$ or $\mathrm{Perv}_{(G_{\jepPubScript{O}})}(\mathrm{Gr},\Bbbk)$, in particular due to the facts that standard/costandard objects have more explicit descriptions and that the ``realization functor''
\[
 D^{\mathrm{b}} \mathrm{Perv}_{\jepPubEuler{IW}}(\mathrm{Gr},\Bbbk) \to D^{\mathrm{b}}_{\jepPubEuler{IW}}(\mathrm{Gr},\Bbbk)
\]
is an equivalence of triangulated categories.

In the analogous setting of Whittaker $\jepPubEuler{D}$-modules over a field of characteristic $0$, this statement already appears in~\cite{jep104ref09}. See Remark~\ref{jep104local026}(2) below for a discussion of possible variants for constructible sheaves over $\jepPubBbb{C}$. Let us also mention the conjecture~\cite[Conj.~59]{jep104ref16} containing this statement as a special case (see~\cite[Ex.~60]{jep104ref16} for more details).

\subsection{Relation with the Finkelberg--Mirkovi\'c conjecture}
\label{jep104local036}

One possible justification for the equivalence~\eqref{jep104local009} comes from a singular analogue of the Finkelberg--Mirkovi\'c conjecture~\cite{jep104ref22}. This conjecture states that, if $\Bbbk$ is a field of positive characteristic $\ell$ bigger than the Coxeter number of $G$, if $I \subset G_{\jepPubScript{O}}$ is an Iwahori subgroup and $I_\mathrm{u} \subset I$ is its pro-unipotent radical, there should exist an equivalence of abelian categories
\[
 \mathsf{F}: \mathrm{Perv}_{I_\mathrm{u}}(\mathrm{Gr},\Bbbk) \xrightarrow{\sim} \mathrm{Rep}_0(G^\vee_\Bbbk)
\]
between the category of $I_\mathrm{u}$-equivariant $\Bbbk$-perverse sheaves on $\mathrm{Gr}$ and the ``extended principal block'' $\mathrm{Rep}_0(G^\vee_\Bbbk)$ of $\mathrm{Rep}(G^\vee_\Bbbk)$, i.e.~the subcategory consisting of modules over which the Harish--Chandra center of the enveloping algebra of the Lie algebra of $G^\vee_\Bbbk$ acts with generalized character $0$.

This equivalence is expected to be compatible with the geometric Satake equivalence in the sense that for $\jepPubEuler{F}$ in $\mathrm{Perv}_{I_\mathrm{u}}(\mathrm{Gr},\Bbbk)$ and $\jepPubEuler{G}$ in $\mathrm{Perv}_{G_{\jepPubScript{O}}}(\mathrm{Gr},\Bbbk)$ we expect a canonical isomorphism
\[
 \mathsf{F}(\jepPubEuler{F} \star^{G_{\jepPubScript{O}}} \jepPubEuler{G}) \cong \mathsf{F}(\jepPubEuler{F}) \otimes \jepPubEuler{S}(\jepPubEuler{G})^{(1)}.
\]
(Here $(-) \star^{G_{\jepPubScript{O}}} (-)$ is the natural convolution action of $\mathrm{Perv}_{G_{\jepPubScript{O}}}(\mathrm{Gr},\Bbbk)$ on the category $\mathrm{Perv}_{I_\mathrm{u}}(\mathrm{Gr},\Bbbk)$, and $(-)^{(1)}$ is the Frobenius twist.)

One might expect similar descriptions for some singular ``extended blocks'' of $\mathrm{Rep}(G^\vee_\Bbbk)$, namely those attached to weights in the closure of the fundamental alcove belonging only to walls parametrized by (non-affine) simple roots, involving some Whittaker-type perverse sheaves.\footnote{This extension of the Finkelberg--Mirkovi\'c conjecture stems from discussions of the fourth author with P. Achar. ``Graded versions'' of such equivalences are established in~\cite{jep104ref01}.} In the ``most singular'' case, this conjecture postulates the existence of an equivalence
\[
 \mathsf{F}_{\mathrm{sing}} : \mathrm{Perv}_{\jepPubEuler{IW}}(\mathrm{Gr},\Bbbk) \xrightarrow{\sim} \mathrm{Rep}_{-\varsigma}(G^\vee_\Bbbk)
\]
between our category of Iwahori--Whittaker perverse sheaves and the extended block of weight $-\varsigma$, where $\varsigma$ is a weight whose pairing with any simple coroot is $1$ (the ``Steinberg block''), which should satisfy
\[
 \mathsf{F}_{\mathrm{sing}}(\jepPubEuler{F} \star^{G_{\jepPubScript{O}}} \jepPubEuler{G}) \cong \mathsf{F}_{\mathrm{sing}}(\jepPubEuler{F}) \otimes \jepPubEuler{S}(\jepPubEuler{G})^{(1)}.
\]
(Here we assume that $\varsigma$ exists, which holds e.g.~if the derived subgroup of $G^{\vee}_\Bbbk$ is simply-connected.)

On the representation-theoretic side, it is well known

that the assignment $V \mapsto \mathsf{L}((\ell-1)\varsigma) \otimes V^{(1)}$ induces an equivalence of categories
\[
 \mathrm{Rep}(G^\vee_\Bbbk) \xrightarrow{\sim} \mathrm{Rep}_{-\varsigma}(G^\vee_\Bbbk),
\]
where $\mathsf{L}((\ell-1)\varsigma)$ is the simple $G^\vee_\Bbbk$-module of highest weight $(\ell-1)\varsigma$; see~\cite[\S II.10.5]{jep104ref29} or~\cite{jep104ref07}.
Our equivalence~\eqref{jep104local009} can be considered a geometric counterpart of this equivalence.

\subsection{Relation with results of Lusztig}

Another hint for the equivalence~\eqref{jep104local009} is given by some results of Lusztig~\cite{jep104ref33}. Namely, in~\cite[\S 6]{jep104ref33} Lusztig defines some submodules $\jepPubEuler{K}$ and $\jepPubEuler{J}$ of (a localization of) the affine Hecke algebra $\jepPubEuler{H}$ attached to $G$. By construction $\jepPubEuler{K}$ is a (non unital) subalgebra of the localization of $\jepPubEuler{H}$, and $\jepPubEuler{J}$ is stable under right multiplication by $\jepPubEuler{K}$. Then~\cite[Cor.~6.8]{jep104ref33} states 

that $\jepPubEuler{J}$ is free as a right based $\jepPubEuler{K}$-module (for some natural bases), with a canonical generator denoted $J_\rho$. Now $\jepPubEuler{H}$ (or rather its specialization at $q=1$) is categorified by the category of Iwahori-equivariant perverse sheaves on the affine flag variety $\mathrm{Fl}$ of $G$. The subalgebra $\jepPubEuler{K}$ (or rather again its specialization) is then categorified by $\mathrm{Perv}_{G_{\jepPubScript{O}}}(\mathrm{Gr},\Bbbk)$ (via the pullback functor to $\mathrm{Fl}$), and similarly $\jepPubEuler{J}$ is categorified by $\mathrm{Perv}_{\jepPubEuler{IW}}(\mathrm{Gr}, \Bbbk)$. From this perspective, the functor in~\eqref{jep104local009} is a categorical incarnation of the map $k \mapsto J_\rho \cdot k$ considered by Lusztig, and the fact that it is an equivalence can be seen as a categorical upgrade of~\cite[Cor.~6.8]{jep104ref33}.

\subsection{Relation with results of Frenkel--Gaitsgory--Kazhdan--Vilonen}

Finally, a third hint for this equivalence can be found in work of the second author with Frenkel, Kazhdan and Vilonen~\cite{jep104ref25,jep104ref26} and more recent work~\cite{jep104ref28}. Working in the context of $\jepPubEuler{D}$-modules over a ground field of characteristic $0$ or $\ell$-adic sheaves over a ground field of arbitrary characteristic, in~\cite{jep104ref26} the authors defined a candidate for the role of a Whittaker category on $\mathrm{Gr}$ using geometry of a complete curve and moduli stacks of bundles over it. A more direct, local definition of such a category is proposed in~\cite{jep104ref28}, where it is also shown that the two constructions produce equivalent categories; the methods of~\cite{jep104ref28} rely on recently developed techniques of $\infty$-categories. Notice also that~\cite[Th.~2.7.1(2)]{jep104ref39} implies an equivalence between the above categories and the Iwahori--Whittaker category.

It was shown in~\cite{jep104ref26} (see in particular~\cite[\S\S 1.2.4--1.2.5]{jep104ref26}) that their Whittaker category is a free right module over the monoidal category $\mathrm{Perv}_{G_{\jepPubScript{O}}}(\mathrm{Gr},\Bbbk)$. Thus, combining these works, we obtain another proof of the equivalence between $\mathrm{Perv}_{\jepPubEuler{IW}}(\mathrm{Gr},\Bbbk)$ and $\mathrm{Perv}_{G_{\jepPubScript{O}}}(\mathrm{Gr},\Bbbk)\cong\mathrm{Rep}(G^\vee_\Bbbk)$, valid when we work with characteristic-$0$ coefficients.

The above results cannot be automatically carried over to our present context, which is that of sheaves with coefficients of positive characteristic. However, the latter equivalence does generalize to our context, and amounts to our equivalence~\eqref{jep104local009}. Of course, we use different methods to prove it.

As explained in~\cite[\S 1.1]{jep104ref26}, in the case of characteristic-$0$ coefficients these properties are closely related to the Casselmann--Shalika formula, and in fact our proof uses the geometric counterpart to this formula known as the \emph{geometric Casselmann--Shalika formula}. (See also~\cite[\S 1.1.1]{jep104ref10} for the relation between the ``Whittaker'' and ``Iwahori--Whittaker'' conditions in the classical setting of modules over the affine Hecke algebra.)

\subsection{Application to tilting objects}

In Section~4 we provide a number of applications of this statement. An important one is concerned with the description of the \emph{tilting objects} in the Satake category. Namely, in the case when $\Bbbk$ is a field of characteristic $\ell$, the \emph{tilting modules} (see e.g.~\cite[Chap.~E]{jep104ref29}) form an interesting family of objects in the category $\mathrm{Rep}(G^\vee_\Bbbk)$. It is a natural question to try to characterize topologically the $G_{\jepPubScript{O}}$-equivariant perverse sheaves on $\mathrm{Gr}$ corresponding to these objects. A first answer to this question was obtained by Juteau--Mautner--Williamson~\cite{jep104ref32}: they showed that, under some explicit conditions on $\ell$, the \emph{parity sheaves} on $\mathrm{Gr}$ for the stratification by $G_{\jepPubScript{O}}$-orbits are perverse, and that their images under~\eqref{jep104local006} are the indecomposable tilting objects in $\mathrm{Rep}(G^\vee_\Bbbk)$. This result was later extended by Mautner and the fourth author~\cite{jep104ref35} to the case when $\ell$ is good for $G$, and it played a crucial role in the proof (by Achar and the fifth author) of the Mirkovi\'c--Vilonen conjecture (or more precisely the corrected version of this conjecture suggested by Juteau~\cite{jep104ref30}) on stalks of standard objects in the Satake category~\cite{jep104ref06}.

It is known (see~\cite{jep104ref32}) that if $\ell$ is bad then the $G_{\jepPubScript{O}}$-constructible parity sheaves on $\mathrm{Gr}$ are not necessarily perverse; so the answer to our question must be different in general. A conjecture was proposed by Juteau--Mautner--Williamson to cover this case, namely that the perverse cohomology objects of the parity complexes are tilting in $\mathrm{Perv}_{G_{\jepPubScript{O}}}(\mathrm{Gr},\Bbbk)$ (so that all tilting objects are obtained by taking direct sums of direct summands of the objects obtained in this way). In our main application we confirm this conjecture, see Theorem~4.10, hence obtain an answer to our question in full generality.

Using this description we prove a geometric analogue of a fundamental result for tilting modules, namely that these objects are preserved by tensor product and by restriction to a Levi subgroup. (On the representation-theoretic side, these results are due to Mathieu~\cite{jep104ref34} in full generality; see~\cite[\S 1.1]{jep104ref32} for more references.) In fact, combined with the Satake equivalence, our proof can also be considered as providing a new complete proof of these properties of tilting modules. In~\cite{jep104ref11}, Baumann and the fourth author also use these facts to obtain a slight simplification of the proof of the geometric Satake equivalence. (Note that the proofs in the present paper do not rely on the latter result.)

\paragraph{Acknowledgements.}

The final stages of this work were accomplished while the fourth author was a fellow of the Freiburg Institute for Advanced Studies, as part of the Research Focus ``Cohomology in Algebraic Geometry and Representation Theory'' led by A. Huber--Klawitter, S. Kebekus and W. Soergel.

We thank P.~Achar and G.~Williamson for useful discussions on the subject of this paper, and the referees for their helpful comments.

